\section{四类之外：按实现平面重看优化栈}
\label{sec:plane}

第 \ref{sec:technique} 节的四类是按「作用于 KV 传输成本的哪个环节」
分的：省量、隐藏、均衡、少步。这个维度覆盖了\textbf{分离本身带来的
额外成本}，但它不是唯一的切法，也不覆盖全部优化。

对照生产系统的实际做法就能看出边界。vLLM 的大规模服务文档列出九项
关键优化，把每 H200 的吞吐从约 1{,}500 提到 2{,}200\,tok/s——其中
\textbf{六项与 KV 传输无关}\texttt{[c-plane]}：异步调度、CUDA graph 的
\texttt{FULL\_AND\_PIECEWISE} 模式、DeepGEMM、DeepEP kernel、
SiLU kernel、算子融合。

这说明四类不够。本节换一个维度重看：按\textbf{实现平面}分。

\subsection{三平面框架}

生产系统普遍把优化分成三个平面\texttt{[c-plane]}：

\begin{table*}[t]
\centering
\caption{三平面与各自的优化手段。与第 \ref{sec:technique} 节的四类
正交——四类只回答「分离的那笔账怎么省」，三平面回答「整个栈有哪几层可以动」。}
\label{tab:plane}
\small
\setlength{\tabcolsep}{9pt}
\begin{tabular}{@{}p{2.5cm}p{5.32cm}p{7.81cm}@{}}
\toprule
平面 & 管什么 & 代表手段 \\
\midrule
数据面 & KV 放哪、留多久 & KV 池分层、打包布局、
  session 感知保留、前缀缓存 \\
执行面 & 单步算多快 & 并行策略（DEP/DCP/PCP/Wide-EP）、
  kernel 融合、CUDA graph、异步调度 \\
控制面 & 谁在什么时候被调度 & 长 prefill token 上限、
  prefill 节拍、路由策略、准入、P/D 配比 \\
\bottomrule
\end{tabular}
\end{table*}

\noindent 这个划分比四类更好用的地方在于：它对应\textbf{不同的负责人
与不同的时间尺度}。数据面是容量规划，执行面是 kernel 与并行团队，
控制面是调度策略。四类都在控制面与执行面的交界处，看不全。

\subsection{执行面：完全不涉及传输的优化}

这是四类漏得最多的一块。它不减少传输、也不隐藏传输，
而是\textbf{直接把单步算得更快}。

\noindent\textbf{算子融合。} 把多个小算子合成一个 kernel，省掉中间的
读写与启动开销。实测幅度不小：

\begin{itemize}[leftmargin=1.3em,itemsep=1pt,topsep=2pt]
  \item SiLU-mul-quant 融合、Cutlass QKV kernel 等改进，与 DBO 一起
        构成 1.5k$\to$2.2k\,tok/s/H200 的主要来源\texttt{[c-plane]}。
  \item Kimi K3 的 GEMM 与 reduce-scatter 融合，改善序列并行的通信
        \texttt{[c-plane]}；latent-tail MoE 融合把端到端延迟降约
        $5\%$\texttt{[c-plane]}。
  \item 用对称内存把 query 多播、部分输出与 LSE 直接写进对端接收槽，
        再在 kernel 内融合——每层延迟降约 $13\%$\texttt{[c-plane]}。
  \item MiniMax M3 的长上下文 indexer（CuteDSL），GB300 上延迟改善
        约 $3\%$ 至 $31\%$，取决于 shape\texttt{[c-plane]}。
\end{itemize}

\noindent 这些数字有个共同特征：\textbf{个位数到三十个百分点，
但都是端到端可见的}。它们不依赖架构选择，改动局部，
是投入产出比最高的一类——也正因如此，它们最容易被
「这是 kernel 团队的事」推给别人，然后从架构讨论里消失。

\noindent\textbf{并行策略：不是「均衡」，是「选哪种切法」。}

我在四类里把负载均衡与并行放在一起谈，这是不准确的。并行策略
是\textbf{选切分维度}，与均衡是两个问题。

\noindent 同样的模型，不同切法的差别是数量级的：
DeepSeek V4 在 32K prompt 上，prefill 上下文并行（PCP8）相对张量并行
（TP8）有 $2.65\times$ 的 prefill 加速\texttt{[c-plane]}。Kimi K3 的
decode 上下文并行（DCP）沿序列切 KV，每个 rank 只留 $1/N$ 的状态，
既降 decode 延迟又扩大 KV 容量\texttt{[c-plane]}。

\noindent 但\textbf{切法必须服从模型架构}，这是本节最重要的一条，
而它是从一个反例得出的。DeepSeek V4 的压缩稀疏注意力包含 indexer、
compressor 与主注意力三部分，上下文并行必须切分并协调全部子层。
vLLM 团队「在重叠通信与优化对应 kernel 上投入很大」，最后
DCP \textbf{只是打平 DEP，没有超过}\texttt{[c-plane]}。

\noindent 原文的总结值得直接引：\textbf{并行必须跟随模型架构；
在一个隐注意力模型上成功的策略，未必能推广到另一个}
\texttt{[c-plane]}。这解释了为什么第 \ref{sec:model} 节按模型给做法
是对的——但理由不是「模型大小不同」，而是\textbf{注意力结构的内部
组成不同，可切的维度就不同}。

\noindent\textbf{第二个反例：流水线并行不适合热轮次。} PP（含 chunked PP）
在长而新的 prompt 上表现很好——大 prefill 提供的计算足以填满各级流水线，
吞吐近线性扩展。但 agentic 场景的多数轮次已有系统提示与前几轮缓存，
每个新请求只增加几百到几千 token，\textbf{没有足够的新计算填满流水线，
bubble 吃掉大部分潜在收益}\texttt{[c-plane]}。

\noindent 原文的措辞很克制：「教训不是 PP 无效，而是它对冷的、
计算重的 prefill 有效，但不该成为热的、前缀重的轮次的默认」
\texttt{[c-plane]}。这是一条\textbf{按负载形态选并行}的判据，
比任何「PP 好/不好」的结论都有用。

\subsection{数据面：容量与命中率}

四类里的 A 类（省传输量）只管「少搬」，但生产系统还有另一半：
\textbf{把 KV 留在更近的地方，提高命中}。

\noindent\textbf{KV 池分层。} 把分布式 KV 池做成多层，用磁盘与纯 CPU
节点扩容量；\texttt{standalone-store} 模式下由外部 Mooncake 客户端
持有 CPU 池与磁盘层，vLLM worker 变成纯请求方\texttt{[c-plane]}。

\noindent\textbf{打包布局。} 这个细节很能说明问题：DeepSeek V4 早期的
KV 布局把不同 cache 类型碎成三个尺寸桶、分配了 \textbf{92 个独立
tensor}，padding 浪费显存，且对 P/D 传输与 offload 都不友好。
新的打包布局把全部 cache group 与层放进每块一个连续分配
\texttt{[c-plane]}，减少描述符与 P/D 传输开销，并允许更小的分配单元
（开 FP4 indexer 时省约 $10\%$ 的 KV 显存）\texttt{[c-plane]}。

\noindent 这一条对本文主题有直接含义：\textbf{数据传输效率不只取决于
网络，还取决于内存布局}。同样的字节数，碎片化布局要发 92 个描述符，
打包布局只发一个。这属于「让传输便宜」，但既不在我的 A 类
（省量）也不在 B 类（隐藏）里。

\noindent\textbf{session 感知的保留。} 对混合了线性/滑窗注意力与全注意力
的模型，前缀复用需要在复用边界上保留线性状态或滑窗缓存。
逐 token 保存太贵，故两条策略并用：按间隔保留每轮 prompt 末尾的状态；
以及「同一前缀被观察到第二次时才保留检查点」\texttt{[c-plane]}。

\noindent 这条的判据是\textbf{观测到第二次才投资}——把存储成本
押在有复用证据的前缀上，而不是预测哪些前缀会被复用。

\subsection{控制面：调度顺序与路由}

\noindent\textbf{调度顺序：一个被四类完全漏掉的维度。}

我原来的 B 类是「让传输不可见」，讲的是 micro-batch 交错。
但控制面还有另一个问题：\textbf{队列里的顺序}。

\noindent vLLM 的 chunked-prefill 调度器默认先进先出。一条长 prefill
会一步接一步地占满整个 token 预算，同 rank 上排队的短轮次
\textbf{在长 prefill 结束前完全无法被调度}——这就是队头阻塞。

\noindent 修法是给单条请求每步可调度的 token 数设上限
（\texttt{--long-prefill-token-threshold}）。用 512 token 的阈值，
长 prefill 就给短轮次留出空间，让它们同批加入、更早开始 decode。
在 B300 上跑 DeepSeek V4 Pro，这使每 GPU 每秒总 token 数（TPGS）
提升最多 $93\%$，P90 交互性约 $2.3\times$\texttt{[c-plane]}。

\noindent 代价也写清了：长请求自身的 TTFT 变高，故对 TTFT 敏感
的部署要用更大的阈值\texttt{[c-plane]}。

\noindent 幅度上，这条比四类里任何一条都大——$93\%$ 与 $2.3\times$，
而 C 类均衡是个位数百分比。原因不难理解：队头阻塞\textbf{让整个
队列白等}，而均衡只抹平尾部差距。

\noindent 还有一条更细的：DEP 下 MoE 的 all-to-all 迫使各 rank 同步推进，
一个 rank 在处理 prefill 就拖慢整组。若 prefill 在各 rank 的不同步
到达，这个代价被反复支付。修法是设 \texttt{--prefill-schedule-interval}，
只每隔 $N$ 个引擎步接纳 prefill，把 prefill 集中到相同的步上，
使其余步成为纯 decode\texttt{[c-plane]}。

\noindent\textbf{路由：一个必须写进正文的反例。}

\noindent 四类的 C 类假设均衡有价值。生产实测推翻了它。

\noindent 观察到聚合部署里各 rank 的 KV 使用严重不均，自然的反应是按
队列深度、运行 token 数或当前 KV 利用率来均衡请求。
\textbf{但这些策略全部劣于简单的 session 粘性路由}\texttt{[c-plane]}。

\noindent 机制值得逐字理解：agentic 会话的轮次间隔很短，下一轮到达时
它的前缀往往还在上一张卡上。把会话迁到更空的 rank，就必须去取 KV——
即使前缀保存在分布式池里。传输是异步的、与计算重叠的，
\textbf{但它不是免费的}：预取的块会临时占用目标 rank 的 GPU KV 容量，
减少它能接纳的序列数。结果是\textbf{队列更均衡了，总并发却更少}
\texttt{[c-plane]}。

\noindent 原文的结论是：轮次间隔短时，\textbf{保住会话局部性比完美
均衡瞬时负载更有价值}；路由必须考虑每个 worker 上已驻留的状态，
而不只是排队的量\texttt{[c-plane]}。

\noindent 这一条对本文的判据是一个补充而不是否定：第
\ref{sec:boundary} 节说分离要考虑负载是否落在窗口内，
现在要加一句——\textbf{即使落在窗口内，均衡策略也可能让收益变薄}，
因为均衡与局部性冲突。

\subsection{四类与三平面是什么关系}

写到这里可以说清两套划分的关系，这也是本节最想给出的结论：

\begin{itemize}[leftmargin=1.3em,itemsep=1.5pt,topsep=2pt]
  \item \textbf{四类是按成本链切的}，覆盖「KV 传输这笔账怎么省」。
        它的优点是能直接对应到收益与代价，缺点是看不到传输之外的层。
  \item \textbf{三平面是按实现层切的}，覆盖「整个栈有哪几层可以动」。
        它的优点是完备，缺点是不能直接告诉你「哪个划算」。
  \item \textbf{两者正交，应该叠加使用}：用三平面保证不漏，
        用四类判断传输那一层的取舍。
\end{itemize}

\noindent 若只能记一句，我选这一句：\textbf{四类解决的是分离引入的
额外成本；而生产系统里最大的单项收益往往来自与分离无关的执行面
改进}——$2.65\times$ 的 PCP、$93\%$ 的调度上限、$13\%$ 的 kernel 融合，
都不在四类里。把注意力全放在「怎么把 KV 搬得更快」，
会错过这些。
